\chapter{Substrates and the runtime}\label{ch:runtime}

\section{The worker}

\texttt{vapor-worker} is a static Zig executable with no libc. At start-up (\textsc{hello}) it fixes the
number of threads; the thread pool and the event counters (\texttt{perf\_event\_open}: cycles,
instructions and cache events when a PMU exists; task clock, page faults and context switches always)
are created \emph{before} the seccomp filter, which is installed with \texttt{TSYNC} on every thread.
Each call carries a \emph{partition descriptor} (a count argument, a grain, pointers that advance per
unit, per-thread scratch, an optional guard) and the worker splits the units among threads. Because the
units are independent rows, the result is the same for one thread or many.

Code runs in one of two modes. \textbf{Native}: the machine code is copied to a writable page, the page
is made read--execute (W\^{}X), the instruction cache is synchronised where the architecture needs it,
and the kernel is called as \texttt{void k(const uint64\_t *args)}. \textbf{Emulated}: RV64GCV code is
interpreted by \texttt{rvemu.zig}, which checks every memory access against the bound buffers and every
fetch against the code, and turns an unknown instruction into an error with its address and word. The
interpreter implements RVV~1.0 faithfully (round-to-nearest, fused \texttt{vfmacc}, canonical NaN,
NaN-boxing, group alignment, \texttt{vill}), with a \emph{poison} mode that writes ones into
tail- and mask-agnostic elements, proving that kernels do not depend on them, and a configurable VLEN
from 128 to 512 bits.

Containment is a seccomp-BPF allow-list (read, write, openat, close, lseek, mmap, munmap, mprotect,
clock\_gettime, setitimer, signals, exit; the architecture is audited),
\texttt{PR\_SET\_NO\_NEW\_PRIVS}, and a \texttt{SIGALRM} watchdog for native code that does not finish.
The tests provoke all four failure classes (\texttt{SIGILL}, \texttt{SIGSEGV}, \texttt{SIGALRM},
\texttt{SIGSYS}); in each, only the worker dies, the BEAM process that owns it reports the signal,
restarts it, and the next unit runs.

\section{The Vulkan daemon}

\texttt{vapor-fabric} drives Vulkan compute through hand-written bindings (libc only for the loader's
\texttt{dlopen}): instance, physical device with a compute queue, logical device, buffers, shader modules
from the SPIR-V emitted by the BEAM, pipelines with push constants, descriptor sets, one command buffer
per run with barriers, per-iteration windows and state copies, submission with a fenced deadline, and
readback. A driver crash or a lost device kills only the daemon. Its plans are self-contained, so a
restarted daemon needs no replay.

\section{One protocol}

Both processes speak the same framing: \texttt{\{packet, 4\}} frames on standard I/O (big-endian
length, little-endian fields), so a process that dies in the middle of a frame cannot desynchronise the
BEAM.

\begin{table}[h]
\centering\small
\begin{tabularx}{\textwidth}{@{}llX@{}}
\toprule
frame & code & content \\
\midrule
\textsc{hello} & 1 & flags, threads $\to$ architecture, sandbox state, version, threads \\
\textsc{run} & 2 & mode, VLEN, flags, fuel, deadline, code, buffers (inline or mapped files), calls with arguments and partition descriptors, iterations, streamed emissions, copies, returns \\
\textsc{emit} & 3 & per-iteration output, streamed \\
\textsc{done} & 4 & elapsed time, retired instructions, counters, returned buffers \\
\textsc{err} & 5 & code, program counter, instruction word, message \\
\textsc{open}, \textsc{step}, \textsc{close} & 6--8 & resident sessions: weights mapped once, state kept between steps, each step writing only its inputs and returning only the requested outputs, with optional state copies inside the worker \\
\bottomrule
\end{tabularx}
\caption{The worker protocol. Tables are sized by the frame itself, so a hostile count cannot allocate
more than the frame describes.}
\end{table}

\section{Memory}

Weights of 64\,KiB or more are written once to \texttt{/dev/shm/vapor-<sha256>} (content addressed,
atomic write) and mapped copy-on-write by every worker; the Vulkan daemon imports them without copy
through \texttt{VK\_EXT\_external\_memory\_host} when the device offers it. Replicas and restarts map the
same files, so $n$ workers serving one model occupy one copy in the page cache. Weights can stay in
bfloat16: the kernels widen sixteen weights exactly on load, so the result is bit for bit that of the
f32 program over the rounded values. Per-token inputs are windows into one buffer; state feedback
($h \leftarrow h_{\text{next}}$) is a copy declared in the plan.

\section{Dispatch and failover}

\texttt{Vapor.Runtime.Dispatch} tries the arbiter's choice and walks down the chain
\[
\text{fabric} \to \text{host AVX-512} \to \text{host base ISA} \to \text{RVV interpreter} \to \text{oracle},
\]
recording the reason for each step. Under the canonical policy every link computes the same bits, so
rerouting never changes an answer.

\section{Parallelism without changing bits}

\begin{table}[h]
\centering\small
\begin{tabularx}{\textwidth}{@{}lXl@{}}
\toprule
form & where & invariant \\
\midrule
SIMD & AVX2, AVX-512, NEON, RVV (VLEN 128--512), SPIR-V & same bits as the oracle \\
instruction-level & $R$ rows per iteration in matrix--vector products, chosen by the allocator & same \\
threads & the worker's pool; decode attention split by key--value head & 1 thread $=$ $n$ threads \\
continuous batching & decode and prefill chunks in one step & a sequence alone $=$ in a batch \\
paged cache & paged key--value writes and attention over a block table & paged $=$ contiguous \\
replicas & one compilation, shared weight pages, shortest queue & same answer from any replica \\
speculation & a draft proposes, the target verifies $k{+}1$ rows in one step & output $=$ the target's \\
GPU & SPIR-V for every model operator & fabric $=$ oracle \\
\bottomrule
\end{tabularx}
\caption{Every form of parallelism draws on independent outputs.}
\end{table}

Speculative decoding needs no further correctness argument: the $k{+}1$-row verification step is
bit-identical to $k{+}1$ single steps by batch invariance, so the output equals the target's greedy
decoding whatever the draft, and the draft only changes the number of steps. Block drafters and tree
drafts are degenerate and general cases of the same rule.

\section{The engine and serving}

\texttt{Vapor.Engine} serves compiled models in resident sessions with continuous batching, a paged
key--value cache, replicas and cancellation: the engine monitors the process that asked for each
sequence, and a sequence whose requester dies leaves the batch at the next step boundary without
changing the bits of the others. Recurrent models (Mamba, Mamba-2, the delta-rule hybrid of
Chapter~\ref{ch:models}) run as step programs whose state stays in the worker. The HTTP server speaks
the OpenAI chat and completion protocols with streaming, and every response carries a receipt, the
canonical digest of model, prompt tokens, parameters and generated tokens. Constrained decoding
(\texttt{Vapor.Grammar}) restricts generation to a grammar derived from a JSON Schema, a regular
expression or a tool's declared arguments, with the invariant that no accepted prefix is a dead end.
